\def\tituloconcepto{QAOA}
% Si el capitulo se compila solo (sin main.tex), esta cabecera abre el documento.
% Antes de % Si el capitulo se compila solo (sin main.tex), esta cabecera abre el documento.
% Antes de % Si el capitulo se compila solo (sin main.tex), esta cabecera abre el documento.
% Antes de \input{cabecera} cada capitulo define \tituloconcepto.
\ifdefined\enmain\else
  \documentclass[11pt,a4paper]{article}
  \input{preambulo}
  \begin{document}
  \title{\tituloconcepto}
  \author{\autor}
  \date{\today}
  \maketitle
\fi
 cada capitulo define \tituloconcepto.
\ifdefined\enmain\else
  \documentclass[11pt,a4paper]{article}
  % Preambulo compartido por todos los documentos. Compilar con XeLaTeX/LuaLaTeX
% (o Tectonic). Si se usa Overleaf, elegir el compilador XeLaTeX.
\usepackage[spanish,es-noshorthands]{babel}
\usepackage{amsmath,amssymb}
\usepackage{booktabs,array}
\usepackage[a4paper,margin=2.5cm]{geometry}
\usepackage{xcolor}
\usepackage{listings}
\usepackage[most]{tcolorbox}
\usepackage{quantikz}
\usepackage[hidelinks]{hyperref}

% ---- Autor (editar a gusto) -------------------------------------------------
\newcommand{\autor}{Genesis Silesky}

% ---- Notacion de Dirac y atajos ---------------------------------------------
% quantikz ya define \ket y \bra; DeclareRobustCommand los redefine sin error.
\DeclareRobustCommand{\ket}[1]{\left|#1\right\rangle}
\DeclareRobustCommand{\bra}[1]{\left\langle#1\right|}
\DeclareRobustCommand{\braket}[2]{\left\langle#1\middle|#2\right\rangle}
\newcommand{\C}{\mathbb{C}}
\newcommand{\R}{\mathbb{R}}
\DeclareMathOperator{\Tr}{Tr}

% ---- Cajas de apoyo --------------------------------------------------------
% "concepto": definiciones previas que hacen falta para entender el capitulo.
\newtcolorbox{concepto}[1]{breakable,colback=blue!4,colframe=blue!55!black,
  fonttitle=\bfseries,title={Concepto previo: #1}}
% "nota": aclaraciones y advertencias.
\newtcolorbox{nota}[1]{breakable,colback=orange!6,colframe=orange!70!black,
  fonttitle=\bfseries,title={#1}}

% ---- Estilos de listados ----------------------------------------------------
\lstdefinestyle{py}{language=Python,basicstyle=\ttfamily\footnotesize,
  breaklines=true,postbreak=\mbox{\textcolor{gray}{$\hookrightarrow$}\space},
  numbers=left,numberstyle=\tiny\color{gray},frame=single,
  commentstyle=\color{green!40!black}\itshape,
  keywordstyle=\color{blue!70!black}\bfseries,
  stringstyle=\color{red!60!black},showstringspaces=false,tabsize=4,
  columns=fullflexible,keepspaces=true}
\lstdefinestyle{salida}{basicstyle=\ttfamily\footnotesize,breaklines=true,
  frame=leftline,backgroundcolor=\color{gray!8},columns=fullflexible,
  keepspaces=true}

  \begin{document}
  \title{\tituloconcepto}
  \author{\autor}
  \date{\today}
  \maketitle
\fi
 cada capitulo define \tituloconcepto.
\ifdefined\enmain\else
  \documentclass[11pt,a4paper]{article}
  % Preambulo compartido por todos los documentos. Compilar con XeLaTeX/LuaLaTeX
% (o Tectonic). Si se usa Overleaf, elegir el compilador XeLaTeX.
\usepackage[spanish,es-noshorthands]{babel}
\usepackage{amsmath,amssymb}
\usepackage{booktabs,array}
\usepackage[a4paper,margin=2.5cm]{geometry}
\usepackage{xcolor}
\usepackage{listings}
\usepackage[most]{tcolorbox}
\usepackage{quantikz}
\usepackage[hidelinks]{hyperref}

% ---- Autor (editar a gusto) -------------------------------------------------
\newcommand{\autor}{Genesis Silesky}

% ---- Notacion de Dirac y atajos ---------------------------------------------
% quantikz ya define \ket y \bra; DeclareRobustCommand los redefine sin error.
\DeclareRobustCommand{\ket}[1]{\left|#1\right\rangle}
\DeclareRobustCommand{\bra}[1]{\left\langle#1\right|}
\DeclareRobustCommand{\braket}[2]{\left\langle#1\middle|#2\right\rangle}
\newcommand{\C}{\mathbb{C}}
\newcommand{\R}{\mathbb{R}}
\DeclareMathOperator{\Tr}{Tr}

% ---- Cajas de apoyo --------------------------------------------------------
% "concepto": definiciones previas que hacen falta para entender el capitulo.
\newtcolorbox{concepto}[1]{breakable,colback=blue!4,colframe=blue!55!black,
  fonttitle=\bfseries,title={Concepto previo: #1}}
% "nota": aclaraciones y advertencias.
\newtcolorbox{nota}[1]{breakable,colback=orange!6,colframe=orange!70!black,
  fonttitle=\bfseries,title={#1}}

% ---- Estilos de listados ----------------------------------------------------
\lstdefinestyle{py}{language=Python,basicstyle=\ttfamily\footnotesize,
  breaklines=true,postbreak=\mbox{\textcolor{gray}{$\hookrightarrow$}\space},
  numbers=left,numberstyle=\tiny\color{gray},frame=single,
  commentstyle=\color{green!40!black}\itshape,
  keywordstyle=\color{blue!70!black}\bfseries,
  stringstyle=\color{red!60!black},showstringspaces=false,tabsize=4,
  columns=fullflexible,keepspaces=true}
\lstdefinestyle{salida}{basicstyle=\ttfamily\footnotesize,breaklines=true,
  frame=leftline,backgroundcolor=\color{gray!8},columns=fullflexible,
  keepspaces=true}

  \begin{document}
  \title{\tituloconcepto}
  \author{\autor}
  \date{\today}
  \maketitle
\fi


\section{QAOA: algoritmo de optimización cuántica aproximada}

\subsection{Objetivo}
Resolver el problema Max-Cut con un circuito variacional de capas alternadas de
costo y de mezcla, y comprobar cómo influye el número de capas $p$.

\subsection{Conceptos previos}
\begin{concepto}{VQE y rotaciones $ZZ$}
Ver capítulo~19 (VQE, regla del desplazamiento) y capítulo~12 (rotaciones y
acoplamiento $\sigma_z\otimes\sigma_z$). QAOA usa la misma idea híbrida: un
circuito parametrizado y un optimizador clásico.
\end{concepto}

\subsection{Fundamento teórico}
QAOA fue propuesto por Farhi, Goldstone y Gutmann en 2014. Es un algoritmo
híbrido cuántico-clásico para problemas de optimización combinatoria (como el
viajante TSP, la coloración de grafos o Max-Cut), aplicables a logística,
planificación de producción o diseño de circuitos.

\paragraph{Algoritmo adiabático cuántico (QAA).} Si un sistema empieza en el
estado fundamental de un Hamiltoniano $H_I$ y $H(t)$ cambia lo bastante despacio
hasta $H_F$, el sistema permanece en el estado fundamental. Se elige $H_I$ con
un estado fundamental fácil de preparar y $H_F$ de modo que su estado
fundamental codifique la solución. Con $H(t)$ pasando de $H_I$ en $t=0$ a $H_F$
en $t=T$, la evolución se simula con la fórmula de Trotter--Suzuki, alternando pasos pequeños de $H_I$ y de
$H_F$. En la práctica, el ruido y los límites físicos impiden las evoluciones
largas que pide la adiabaticidad.

\paragraph{QAOA.} Adopta la misma idea con evoluciones más cortas y rápidas, y
busca una solución \emph{aproximada} optimizando parámetros. Los pasos son:
\begin{enumerate}
  \item \textbf{Hamiltonianos:} se define un Hamiltoniano de costo $H_C$ (el
    estado de mayor energía codifica la solución) y uno de mezcla $H_M$ que no
    conmuta con $H_C$ (aquí, $H_M=\sum_iX_i$).
  \item \textbf{Estado inicial:} $\ket s=\ket+^{\otimes n}$, el estado de mayor
    energía de $H_M$.
  \item \textbf{Circuito:} la capa de costo $U_C(\gamma)=e^{-i\gamma H_C}$ y la
    de mezcla $U_M(\beta)=e^{-i\beta H_M}$ (que es $R_x(2\beta)$ en cada qubit).
  \item \textbf{Evolución:} con $p$ capas, $\gamma=(\gamma_1,\ldots,\gamma_p)$ y
    $\beta=(\beta_1,\ldots,\beta_p)$:
    \begin{equation}
      \ket{\psi_p(\gamma,\beta)}=e^{-i\beta_pH_M}e^{-i\gamma_pH_C}\cdots
      e^{-i\beta_1H_M}e^{-i\gamma_1H_C}\ket+^{\otimes n}.
    \end{equation}
  \item \textbf{Valor esperado:}
    $E(\gamma,\beta)=\langle\psi_p|H_C|\psi_p\rangle$.
  \item \textbf{Optimización:} un método clásico actualiza $\gamma$ y $\beta$
    para \emph{maximizar} $E$.
\end{enumerate}

\paragraph{Max-Cut.} Dado un grafo $G=(V,E)$ con $n$ vértices, se reparten los
vértices en dos grupos para \emph{maximizar} el número de aristas entre grupos
distintos. Si $x_i\in\{0,1\}$ dice a qué grupo (por ejemplo, a qué día) pertenece
el vértice $i$, la función a maximizar es
\begin{equation}
  L(x)=\sum_{(u,v)\in E}\big(x_u+x_v-2x_ux_v\big).
\end{equation}
Para dos lugares turísticos unidos por una arista hay $2^2=4$ formas de
repartirlos: $00$ y $11$ cortan $0$ aristas y $01$ y $10$ cortan $1$. Para pasar
a un Hamiltoniano se sustituye $x_i\to\frac{I-Z_i}2$ (vale $0$ en $\ket0$ y $1$ en
$\ket1$) y queda
\begin{equation}
  H_C=\sum_{(u,v)\in E}\frac{1-Z_uZ_v}{2}.
\end{equation}

\paragraph{Cómo se hace $e^{-i\gamma Z_uZ_v}$.} Con una CNOT, una rotación
$R_z$ sobre el segundo qubit y otra CNOT. En el ejemplo, la capa de costo con
ángulo de puerta $a=-\gamma$ coincide con $e^{-i\gamma H_C}$ salvo una fase
global.

\subsection{Código (completamente comentado)}
\lstinputlisting[style=py]{../../fuentes/20_qaoa.py}

\subsection{Explicación celda por celda}
\begin{description}
  \item[Celda 2] Lista las cuatro formas de repartir los dos lugares y cuántas
    aristas corta cada una.
  \item[Celda 3] Construye $H_C$ y comprueba que su valor máximo es $1$.
  \item[Celda 4] Arma el circuito QAOA y verifica que la capa de costo es
    $e^{-i\gamma H_C}$.
  \item[Celda 5] Calcula $\langle H_C\rangle$ y su gradiente. Como un mismo
    parámetro aparece en varias puertas (la mezcla usa el mismo $\beta$ en todos
    los qubits), la regla del desplazamiento se aplica a cada puerta por separado y se suman
    las contribuciones. Se compara con la derivada numérica.
  \item[Celda 6] Ascenso por gradiente con $p=1$.
  \item[Celda 7] Muestra las probabilidades de cada cadena con los parámetros
    óptimos.
  \item[Celda 8] Repite con $p=1,2,3$, como pide la práctica.
\end{description}

\subsection{Salida obtenida}
\lstinputlisting[style=salida]{../../fuentes/salidas/20_qaoa.txt}
(El dibujo del circuito no se reproduce aquí.) Es determinista (la
optimización usa una semilla fija).

\subsection{Interpretación}
El gradiente por desplazamiento coincide con el numérico
($-0.44849$ y $0.08848$). Con $p=1$ el valor esperado llega a $1.0000$ (el
máximo): la probabilidad se reparte entre $01$ y $10$ ($0.5$ cada una), que son
las dos formas de cortar la arista, y $00$ y $11$ quedan en $0$. Con $p=2$ y
$p=3$ también se obtiene $1.0000$: en este grafo tan pequeño, una sola capa
alcanza.

\subsection{Ejercicios}
\begin{enumerate}
  \item Cambia el grafo a un triángulo (aristas $(0,1),(1,2),(0,2)$, $n=3$).
    ¿Cuál es el valor máximo de $H_C$? ¿Llega QAOA con $p=1$?
  \item Cambia la tasa de aprendizaje y el número de iteraciones y mira cómo
    cambia el resultado.
  \item Prueba otras semillas para los parámetros iniciales con $p=1$: ¿llega
    siempre a $1$?
\end{enumerate}

\subsection{Errores comunes}
\begin{itemize}
  \item Confundir el signo: QAOA aquí \emph{maximiza} $\langle H_C\rangle$ (el
    número de aristas cortadas), mientras que VQE minimiza una energía.
  \item Olvidar que los parámetros compartidos necesitan sumar las
    contribuciones de cada puerta al calcular el gradiente.
\end{itemize}

\ifdefined\enmain\else
  \end{document}
\fi

